\scanpage{95}\mspage{39}
% L1-3
Il est difficile d'imaginer que $\bar k$, ou ce qui revient au $=$ $\bar{\unsure{\varphi|\psi}}$, puisse être une équivalence, sans qu'il existe au moins une ``$W_{\unsure{f|S}}$-homotopie''
% L4
\diagram{\begin{tikzcd}[ampersand replacement=\&,column sep=0.43em]
 \deleted{\unsure{\text{entre}}}\ \deleted{\unsure{\varphi k}}\quad
 \mathrm{id}_{\underline{\mathrm{Fibsc}}_{\deleted{\unsure{?}}0}(\deleted{\unsure{E}}S)}
 \arrow[r]\arrow[rrrr,phantom,"h"]\&{}\&{}\arrow[l]\arrow[r]\&{}\&{}\arrow[l]\arrow[r]\&{}\&{}\arrow[l]\arrow[r]\&\deleted{\unsure{\bar{\phantom{k}}}}\varphi k
\end{tikzcd}}
% L5-7
dans la catégorie des endofoncteurs de $\underline{\mathrm{Fibsc}}_0\,S$. \deleted{Donc il faut, pour $X$ objet donné dans cette catégorie,}

% L7-12
À vrai dire, je me rends compte maintenant que la description de $\bar\varphi$ à partir de $\varphi:X\longmapsto\Phi_0 X=\Phi C(X)$ est inutilement compliquée, car on avait fait le détour par $C(X)$ uniquement pour être sûr de tomber dans $\underline{\mathrm{Fibsc}}_0\,S$, pas seulement dans $\underline{\mathrm{Fibsc}}$.
% L13-17
Mais si on laisse tomber $\underline{\mathrm{Fibsc}}_0\,S$, on s'en fout, et le foncteur $\unsure{\varphi'|\varphi}:X\mapsto\Phi(X)$ est beaucoup plus naturel. Et \deleted{\unsure{si}} si on ne laisse pas tomber, il n'y a pas d'inconvénient de remplacer la source $\underline{\mathrm{Fib}}\,S$
% L18-22
par $\underline{\mathrm{Fib}}_0\,S$ (lequel a $=$ localisé, à équi. près), et sur $\underline{\mathrm{Fib}}_0\,S$ à nouveau il n'y a pas lieu de s'encombrer de $CX$, pour le plaisir de lui appliquer $\Phi$, car $\Phi(X)$ lui-même est déjà dans $\underline{\mathrm{Fibsc}}_0\,S$. Donc, la question est celle-ci : Si $X$ varie
% archive sheet number: 94
